\section*{ЗАКЛЮЧЕНИЕ}
\addcontentsline{toc}{section}{ЗАКЛЮЧЕНИЕ}

Итогом выпускной квалификационной работы стала онлайн-платформа с базой данных, предназначенная для поиска работы и подбора IT-специалистов. Основное внимание при разработке уделялось базовому циклу взаимодействия между соискателем и работодателем: публикации вакансий, поиску предложений, отправке откликов и обработке заявок.

Первая глава посвящена предметной области IT-рекрутинга. В ней отмечены особенности подбора IT-специалистов: значимость специализации, конкретных навыков, формата занятости и прозрачности статуса отклика. В этом же разделе выделены основные роли пользователей: гость, соискатель, работодатель и администратор. Проведённый анализ показал, что для решения выявленных проблем требуется единая web-платформа, объединяющая вакансии, профили кандидатов, отклики и средства административного контроля.

В техническом задании закреплены требования к системе: состав хранимых данных, функции для каждой роли, требования к интерфейсу, надёжности, информационной безопасности и совместимости. Варианты использования, описанные в этом разделе, охватывают основные сценарии взаимодействия пользователей с платформой.

Технический проект определил архитектуру приложения. Серверная часть построена на Node.js и Express, данные хранятся в PostgreSQL, для работы с базой данных используется Sequelize ORM. Клиентская часть реализована на HTML, CSS и JavaScript. Также спроектирована схема базы данных, в которую вошли таблицы пользователей, профилей соискателей, компаний, вакансий, откликов, избранных вакансий и навыков.

В процессе реализации созданы ключевые программные компоненты: регистрация и авторизация, профиль соискателя, компании работодателей, публикация и фильтрация вакансий, отклики, избранное, навыки, административное управление. Доступ к функциям разграничен по ролям: каждый пользователь получает только те возможности, которые соответствуют его роли.

Платформа реализует основные сценарии для каждой роли. Соискатель регистрируется, заполняет профиль с профессиональными сведениями, ищет вакансии, откликается и добавляет предложения в избранное. Работодатель создаёт компанию, публикует вакансии, просматривает отклики, меняет их статус и ищет кандидатов. Администратор контролирует пользователей и вакансии: может блокировать учётные записи и удалять некорректные предложения.

Системное тестирование охватило основные сценарии: регистрацию, вход, работу с профилем соискателя, поиск и фильтрацию вакансий, отклики, избранное, создание компании, публикацию вакансии, просмотр и изменение статусов откликов, поиск кандидатов, администрирование и проверку прав доступа. Результаты тестирования показали корректную работу проверенных функций.

Отдельно проверена обработка ошибок и ограничений. Без авторизации закрытые разделы недоступны, действия ограничены ролью пользователя, повторный отклик на одну вакансию или дублирование вакансии в избранном не выполняются. Это подтверждает выполнение требований по надёжности и информационной безопасности, заложенных в техническом задании.

Поставленная цель ВКР достигнута. Создана web-платформа, хранящая и обрабатывающая данные пользователей, профилей, компаний, вакансий, откликов, избранного и навыков, а также поддерживающая ключевые процессы поиска работы и подбора IT-специалистов.

В дальнейшем систему можно развивать. Среди возможных направлений: уведомления, восстановление пароля, загрузка PDF-резюме, расширенная модерация, личные сообщения, а также рекомендательные механизмы для подбора вакансий и кандидатов. 
